% EMNLP 2025 Paper Draft
% Title: Same Signal, Opposite Meaning: Why Adaptive Compute Fails Across Environments
% Template: EMNLP (ACL style)

\documentclass[11pt]{article}

% === EMNLP / ACL Style ===
% NOTE: Replace with the actual EMNLP 2025/2026 style file when available.
% Download from: https://github.com/acl-org/acl-style-files
\usepackage[hyperref]{emnlp}  % or \usepackage{acl}
\usepackage{times}
\usepackage{latexsym}
\usepackage{amsmath}
\usepackage{amssymb}
\usepackage{amsthm}
\usepackage{graphicx}
\usepackage{booktabs}
\usepackage{multirow}
\usepackage{enumitem}
\usepackage{xcolor}
\usepackage{subcaption}
\usepackage{pifont}

% === Theorem Environments ===
\newtheorem{proposition}{Proposition}
\newtheorem{definition}{Definition}

% === Shortcuts ===
\newcommand{\cmark}{\ding{51}}
\newcommand{\xmark}{\ding{55}}

\title{Same Signal, Opposite Meaning:\\Why Adaptive Compute Fails Across Environments}

% Anonymous submission for review
\author{Anonymous}

\begin{document}
\maketitle

% ============================================================
% ABSTRACT
% ============================================================
\begin{abstract}
Adaptive test-time compute methods for LLM agents---rollout evaluation,
tree search, $K$-variant sampling---improve decision quality but incur
$5$--$15\times$ overhead. A growing family of \emph{selective triggering}
methods reduces this cost by using uncertainty signals (token entropy,
calibrated confidence, vote disagreement) to decide when extra compute
is worthwhile, implicitly assuming a \emph{fixed} mapping from signal
value to compute need. We show that this assumption is wrong. Through
systematic measurement across 8 diverse agent environments, we find
that the signal--utility direction \emph{reverses}: token entropy
correlates negatively with rollout utility in fact-verification
($\rho{=}{-}0.119$, FEVER) but positively in code generation
($\rho{=}{+}0.317$, APPS Interview), while showing no statistically
significant correlation in introductory coding
($\rho{=}{+}0.012$, $p{=}0.63$, APPS Intro).
The \emph{identity} of the most informative signal also varies
entirely---\texttt{step\_count} dominates in HotpotQA
($\rho{=}{-}0.494$) while \texttt{num\_available\_actions} dominates
in WebShop ($\rho{=}{+}0.444$). We explain this reversal via a
\emph{two-source uncertainty model}: the same entropy reflects either
information poverty (rollouts cannot help) or decision difficulty
(rollouts explore viable alternatives), an instance of Simpson's
paradox in the signal--utility space. We prove that direction discovery
is a necessary condition for cross-environment non-negative value of
computation. Building on this analysis, we propose EAAG
(Environment-Aware Adaptive Gating), which lets the agent
autonomously discover environment-specific gating patterns through
exploration, LLM-based reasoning, and LASSO-based direction learning,
with zero per-step deployment cost. Across 8 evaluation environments,
EAAG achieves 34 wins (18 statistically significant at 95\% CI)
vs.\ 2 losses against 6 competing methods
in head-to-head SR comparisons and exhibits environment-adaptive
gating patterns: trigger rate varies with environment---73\% in rollout-safe
TWExpress, 66\% in high-headroom HotpotQA, 28\% in WebShop---without
explicit headroom estimation.
\end{abstract}

% ============================================================
% 1. INTRODUCTION
% ============================================================
\section{Introduction}
\label{sec:intro}

% P1: Background + existing methods
Large language model (LLM) agents benefit substantially from
test-time compute---rollouts, verification, and tree search---but
these optimizers incur 5--15$\times$ computational
overhead~\citep{yao2023tree, hao2023reasoning, zhou2024lats}.
A growing body of work addresses this inefficiency through
\emph{selective triggering}: using uncertainty signals to decide
\emph{when} extra compute is worthwhile, rather than applying it
uniformly~\citep{catts2026, seag2025, cats2025, corefine2026,
adaptthink2025, thinkless2025}. Across diverse mechanisms---vote
entropy~\citep{catts2026}, calibrated
confidence~\citep{cats2025, seag2025}, token-level
entropy~\citep{corefine2026, thinkjustenough2025}, and reinforcement
learning~\citep{adaptthink2025, thinkless2025}---each method has
demonstrated effectiveness within its respective evaluation setting.

% P2: Hidden assumption
Despite differing mechanisms---vote-derived
uncertainty~\citep{catts2026}, entropy
thresholding~\citep{corefine2026, thinkjustenough2025},
calibrated confidence~\citep{cats2025, seag2025}, and reinforcement
learning~\citep{adaptthink2025, thinkless2025}---these methods share
an unquestioned assumption: \emph{the direction of the
signal--utility correlation is fixed}. High entropy means the agent
is struggling and needs help; low confidence means a rollout would
be valuable. This mapping from signal to compute need is treated as
a given, never explicitly verified.

% P3: The assumption is wrong + consequences
We show that this assumption is wrong. Through systematic measurement
across 8 diverse agent environments, we find that the signal--utility
direction \emph{reverses}: token entropy correlates negatively with
rollout utility in fact-verification tasks
($\rho{=}{-}0.119$, FEVER) but positively in code generation
($\rho{=}{+}0.317$, APPS Interview), while showing no statistically
significant correlation in introductory coding
($\rho{=}{+}0.012$, $p{=}0.63$, APPS Intro). Beyond direction, the \emph{identity} of the most informative
signal varies entirely across environments---\texttt{step\_count}
dominates in HotpotQA ($\rho{=}{-}0.494$) while
\texttt{num\_available\_actions} dominates in WebShop
($\rho{=}{+}0.444$); single-signal entropy achieves AUC${\approx}$0.50
(indistinguishable from chance) while multi-signal gates reach
AUC${\approx}$0.83.
All fixed-direction methods systematically fail
in at least two of eight evaluation environments; on FEVER, CATTS
achieves 34.2\%---\emph{below} the 37.0\% no-trigger baseline.
The cost of wrong direction is catastrophic: reversing the discovered
direction causes SR to drop by 37.0\,pp on HotpotQA,
with MLP falling to 45.3\%---below the 49.0\% no-trigger baseline.
When the direction is wrong, more precise calibration makes
performance \emph{worse}, not better.

% P4: Why + problem definition
We explain this reversal via a \emph{two-source uncertainty model}
(\S\ref{sec:toy-model}): the same entropy signal reflects two
fundamentally different sources---\emph{information poverty} (where
the agent lacks data and rollouts cannot help) and \emph{decision
difficulty} (where the agent faces multiple viable paths and rollouts
explore them). Environments differ in their proportion of these two
state types, causing the same signal to carry opposite meaning. This
is an instance of Simpson's paradox~\citep{simpson1951interpretation,
pearl2014understanding} in the signal--utility space: aggregating
heterogeneous subpopulations with opposing within-group trends
reverses the aggregate correlation. More broadly, our findings reveal
that \emph{uncertainty signal semantics are not an intrinsic property
of the signal but a function of the environment's information
structure}---a distinction that all prior methods implicitly collapse.
We prove that direction discovery is a \emph{necessary condition} for
cross-environment non-negative value of computation
(Proposition~\ref{prop:necessity}).

% P5: Method + results
We propose EAAG (Environment-Aware Adaptive Gating), which enables
LLM agents to autonomously discover environment-specific compute
allocation patterns. EAAG operates in three steps: (1)~the agent
\emph{explores} a new environment with 50 episodes of randomized
gating; (2)~an LLM \emph{reasons} about exploration data to identify
task-specific patterns and generate feature hypotheses; (3)~LASSO
\emph{learns} signal direction and importance, training a logistic
gate with zero per-step deployment cost. Unlike all prior methods,
EAAG requires no Phase~1 calibration data and no human-specified
signal direction. Across 8 evaluation environments, EAAG
achieves 34 wins (18 statistically significant at 95\% CI)
vs.\ 2 losses against 6 competing methods in
head-to-head SR comparisons, Pareto-dominating in 5/6 shared
environments. The
learned gate exhibits environment-adaptive gating patterns: trigger rate
varies across environments---73\% in rollout-safe TWExpress, 66\%
in high-headroom HotpotQA, 28\% in WebShop---without explicit
headroom estimation.

% P6: Contributions
Our contributions are:
\begin{enumerate}[leftmargin=*,nosep]
\item \textbf{Finding $+$ theory.} We discover that signal--utility
  direction reverses across environments (8 environments, 4
  observations) and explain this via a two-source uncertainty model
  grounded in Simpson's paradox and the epistemic/aleatoric
  distinction. We prove that direction discovery is a necessary
  condition for cross-environment generalization
  (Proposition~\ref{prop:necessity}).

\item \textbf{EAAG method.} We propose Environment-Aware Adaptive
  Gating, where the LLM autonomously discovers environment-specific
  gating patterns via exploration, reasoning, and LASSO-based
  direction learning. Zero per-step overhead, no calibration data
  required.

\item \textbf{Systematic evaluation.} EAAG achieves 34 wins
  (18 statistically significant) vs.\
  2 losses across 8 diverse environments with
  environment-adaptive gating patterns matching rollout headroom.
\end{enumerate}


% ============================================================
% 2. RELATED WORK
% ============================================================
\section{Related Work}
\label{sec:related}

\subsection{Adaptive Test-Time Compute}
\label{sec:related-adaptive}

\paragraph{Signal-based methods.}
SEAG~\citep{seag2025} triggers search when mean token confidence
falls below a threshold. CaTS~\citep{cats2025} uses Platt-scaled
confidence for calibrated early stopping. CoRefine~\citep{corefine2026}
refines outputs when per-token entropy exceeds a threshold. Think
Just Enough~\citep{thinkjustenough2025} applies entropy-based early
stopping with a fixed threshold. DiffAdapt~\citep{diffadapt2025}
trains a lightweight probe for difficulty estimation, assuming a
universal U-shaped entropy pattern. All assume a fixed relationship
between their chosen signal and compute need.

\paragraph{Vote-based methods.}
CATTS~\citep{catts2026} samples $K{=}5$ candidate actions, computes
vote entropy and margin, and triggers a full evaluation when
disagreement is high. This incurs $K$ additional forward passes per
step regardless of whether the optimizer is invoked.

\paragraph{RL-based methods.}
AdaptThink~\citep{adaptthink2025} and
Thinkless~\citep{thinkless2025} learn think/no-think policies via
reinforcement learning. ARPO~\citep{arpo2025} uses entropy-based
adaptive rollout at tool-call steps. These methods implicitly learn
direction through RL training but differ from signal-based methods in
three fundamental ways: (1)~they require \emph{per-environment RL
training} (thousands of episodes with reward signals), making them
incompatible with zero-shot deployment to new environments;
(2)~the learned policy is a black-box neural network that provides no
interpretable direction signal---practitioners cannot inspect
\emph{why} the agent triggers compute; (3)~they operate in the
\emph{reasoning} setting (single-turn chain-of-thought), not the
\emph{multi-step agent} setting where the optimizer $T$ is an
external process (rollout, tree search) invoked at each decision
step. For these reasons, RL-based methods are not directly comparable
to signal-based gating methods and we do not include them as
baselines.

\medskip\noindent
Despite differing mechanisms, all these methods share a common
assumption: a fixed relationship between uncertainty signals and
compute need. Our finding---that this relationship reverses across
environments---challenges this assumption and motivates
direction-aware gating.

\subsection{Orthogonal Work}
\label{sec:related-orthogonal}

Test-time scaling~\citep{snell2024scaling} studies compute-optimal
allocation at the problem level, finding that effectiveness varies
critically with prompt difficulty; we address the complementary
step-level question and show that signal \emph{meaning}---not just
effectiveness---varies across environments. Search methods such as
LATS~\citep{zhou2024lats} and FLARE~\citep{wang2026flare} improve
the search process itself; we control \emph{when} to search.
Recent work on LLM uncertainty~\citep{tao2025revisiting,
heo2025llmuncertainty} independently observes that calibration
quality varies across task types, providing converging evidence that
uncertainty semantics are domain-dependent.

\textbf{Concurrent work.} Several 2025--2026 papers independently
explore adaptive compute: CATTS~\citep{catts2026} uses vote-derived
gating for web agents; AdaptThink~\citep{adaptthink2025} learns
think/no-think via RL; DiffAdapt~\citep{diffadapt2025} uses a probe
for difficulty estimation. While sharing our motivation for selective
triggering, all assume or implicitly learn a fixed signal direction.
Our work is the first to empirically demonstrate that signal--utility
direction reverses across environments and to propose direction-aware
gating as a solution.


% ============================================================
% 3. SIGNAL-UTILITY LANDSCAPE
% ============================================================
\section{The Signal--Utility Landscape}
\label{sec:signal-landscape}

Before proposing a solution, we characterize the relationship between
uncertainty signals and optimizer utility across diverse environments.
If the fixed-direction assumption of prior methods holds, we should
observe consistent signal--utility correlations; if it fails, the
failure mode will guide method design.

We study LLM-based agents in multi-step interactive environments.
At each step $t$, the agent observes state $s_t$, selects action
$a_t$, and transitions to $s_{t+1}$. An environment-specific
\emph{test-time optimizer} $T$ (e.g., rollout evaluation, $K$-variant
sampling) can be invoked at any step to potentially improve action
quality. We define the \emph{optimizer utility} as:
\begin{equation}
  U(T, s) = \mathbb{E}[R(\tau) \mid a{=}T(s)]
           - \mathbb{E}[R(\tau) \mid a{=}\pi(s)]
  \label{eq:utility}
\end{equation}
where $R(\tau)$ is the episode return. $U > 0$ indicates the optimizer
improves the outcome; $U \leq 0$ indicates it is wasteful or harmful.
The \emph{adaptive gating problem} is to learn a gate
$g: \mathcal{S} \to \{0, 1\}$ that triggers $T$ only when $U > 0$.

\subsection{Empirical Landscape}
\label{sec:empirical-landscape}

\textbf{Observation 1: The signal--utility direction reverses across
environments \emph{and} model backbones.}
Table~\ref{tab:signal-discovery} reports Spearman correlations between
observable signals and optimizer utility across 8 environments and 3
model backbones (Qwen3-4B, Phi-3.5-mini, Llama-3.1-8B).
Token entropy---the signal used by the majority of prior
methods---exhibits $\rho{=}{-}0.119$ in FEVER but $\rho{=}{+}0.317$
in APPS Interview on Qwen3-4B. Moreover, the \emph{same signal in the
same environment} can reverse across backbones: FEVER entropy is
$\rho{=}{-}0.156$ on Phi-3.5 but $\rho{=}{+}0.428$ on Llama-3.1;
HotpotQA entropy is $\rho{=}{-}0.346$ on Llama but $\rho{=}{+}0.184$
on Phi-3.5. Structural signals like \texttt{step\_count} show more
consistency across backbones (all negative in 4 main environments),
while model-dependent signals like entropy are fundamentally
unpredictable. The same numerical signal has opposite meaning depending
on \emph{both} the environment and the model.

\begin{table}[t]
\caption{Signal--utility correlations across 8 environments and 3 backbones.
$\rho$: Spearman correlation of token entropy with optimizer utility $U$.
$^*$: significant at $p{<}0.05$.
Direction varies across \emph{both} environments and model backbones.}
\label{tab:signal-discovery}
\centering\small
\begin{tabular}{lcccc}
\toprule
Environment & Qwen3-4B & Phi-3.5 & Llama-3.1 & Consistent \\
\midrule
HotpotQA       & $-$0.327$^*$ & $+$0.184$^*$ & $-$0.346$^*$ & {\color{red}$\times$} \\
APPS Intro     & $+$0.144$^*$ & $-$0.129$^*$ & $-$0.242$^*$ & {\color{red}$\times$} \\
WebShop        & $+$0.133$^*$ & $+$0.335$^*$ & $+$0.287$^*$ & {\color{green}\checkmark} \\
FEVER          & $-$0.119$^*$ & $-$0.156$^*$ & $+$0.428$^*$ & {\color{red}$\times$} \\
APPS Intv.     & $+$0.317$^*$ & $-$0.024     & $-$0.249$^*$ & {\color{red}$\times$} \\
TWExpress      & $-$0.290$^*$ & $+$0.000     & $+$0.000     & --- \\
CRUXEval       & $-$0.048     & $-$0.065     & $-$0.045     & --- \\
Plancraft      & $-$0.016     & $+$0.000     & $-$0.176$^*$ & --- \\
\bottomrule
\end{tabular}
\end{table}

\textbf{Observation 2: The identity of the most informative signal
varies across environments.} EAAG's LASSO selects different feature
subsets per environment (Figure~\ref{fig:feature-heatmap}):
\texttt{step\_count} is selected in 6/7 environments (most universal),
but WebShop relies on \texttt{num\_available\_actions} and LLM-generated
features (\texttt{price\_mentioned}, \texttt{action\_is\_click}).
No single signal is universally informative.

\textbf{Observation 3: Single scalar signals carry near-zero
information.} AUC analysis across 4 environments with available
probe data (Figure~\ref{fig:auc-hierarchy}) reveals a clear hierarchy:
single token entropy achieves AUC${\approx}$0.50 (indistinguishable
from chance), multi-signal logistic regression achieves
AUC${\approx}$0.83, and hidden-state probes reach AUC${\approx}$0.90. The information needed
to predict rollout value simply does not exist in a single scalar
signal at the per-step level.

\textbf{Observation 4: All fixed-direction methods fail systematically.}
Head-to-head SR comparison across 8 environments yields 34 wins
(18 statistically significant at 95\% bootstrap CI) and
2 losses for EAAG against 6 competing methods
(Table~\ref{tab:winloss}). Every fixed-direction baseline fails in at
least 2 environments. Most strikingly, CATTS achieves 34.2\% on
FEVER---below the 37.0\% no-trigger baseline---demonstrating that
wrong-direction gating is actively harmful.


\subsection{Why Does Direction Reverse? A Two-Source Model}
\label{sec:toy-model}

\paragraph{Intuition.}
Consider a web-shopping agent choosing between products. When it
hesitates (high entropy) because multiple items match the query,
exploring alternatives via rollouts is valuable---the agent has
adequate information but faces a genuine choice. Now consider a
fact-verification agent that hesitates because the retrieved evidence
is ambiguous. Here, high entropy reflects \emph{missing information},
not optionality: rollouts from the same model with the same evidence
will not resolve the ambiguity. In the first case, entropy positively
predicts rollout utility; in the second, it predicts rollout
\emph{futility}. The same numerical signal carries opposite meaning
because the \emph{source} of uncertainty differs. We formalize this
intuition below.

\paragraph{Model.}
We propose a parsimonious model explaining when and why the
signal--utility direction reverses. Consider two sources of
uncertainty at each decision step:

\begin{itemize}[leftmargin=*,nosep]
\item \textbf{Type~I (information poverty):} The agent lacks
  sufficient information. High entropy reflects confusion, not
  optionality. Rollouts from the same model cannot compensate for
  missing information, so utility is \emph{negatively} related to
  entropy: $U_I(s) \sim -\alpha \cdot H(s) + \varepsilon_I$,
  \ $\alpha > 0$.

\item \textbf{Type~D (decision difficulty):} The agent has adequate
  information but faces multiple viable paths. High entropy reflects
  optionality. Rollouts can exploit this diversity, so utility is
  \emph{positively} related to entropy:
  $U_D(s) \sim +\beta \cdot H(s) + \varepsilon_D$, \ $\beta > 0$.
\end{itemize}

Let $p_I(\mathcal{E})$ denote the fraction of Type~I states in
environment $\mathcal{E}$. The marginal correlation becomes:
\begin{equation}
  \rho(\mathcal{E}) \;\approx\;
  \beta - (\alpha + \beta)\,p_I(\mathcal{E})
  \label{eq:direction}
\end{equation}
Direction reversal occurs at the critical threshold
$p_I^* = \beta/(\alpha+\beta)$: environments with $p_I > p_I^*$
exhibit negative $\rho$ (Type~I dominated), while those with
$p_I < p_I^*$ exhibit positive $\rho$ (Type~D dominated). Near
$p_I^*$, the two effects cancel and entropy carries negligible
signal ($\rho \approx 0$).

\paragraph{Environment mapping.}
Table~\ref{tab:env-type-mapping} maps each evaluation environment to
its dominant uncertainty type based on task structure.
This mapping is primarily determined by the environment's information
structure, but multi-backbone experiments
(Appendix~\ref{app:multi-backbone}) reveal that the model backbone
modulates the effective $p_I$: different models may process the same
environment with different balances of Type~I vs Type~D uncertainty,
causing entropy $\rho$ to flip sign (e.g., FEVER: $\rho{=}{-}0.156$
on Phi-3.5 but $\rho{=}{+}0.428$ on Llama-3.1). Structural signals
like \texttt{step\_count} are more robust across backbones.

\begin{table}[t]
\caption{Environment-to-type mapping. The dominant uncertainty
type is shaped by the environment's information structure and
modulated by the model backbone.}
\label{tab:env-type-mapping}
\centering\small
\begin{tabular}{llcc}
\toprule
Environment & Dominant Type & Pred.\ $\rho$ & Obs.\ $\rho$ \\
\midrule
HotpotQA       & Type~I (info-poverty) & $-$ & $-0.041$ \\
FEVER          & Type~I (info-poverty) & $-$ & $-0.119$ \\
TWExpress      & Type~I (info-poverty) & $-$ & $-0.290$ \\
\midrule
APPS Intv.     & Type~D (decision-diff) & $+$ & $+0.317$ \\
\midrule
APPS Intro     & Mixed ($p_I \approx p_I^*$) & $\approx 0$ & $+0.012$ \\
WebShop        & Mixed (entropy weak) & $\approx 0$ & $-0.019$ \\
\midrule
CRUXEval       & Weak signal & $\approx 0$ & $-0.048$ \\
Plancraft      & Weak signal & $\approx 0$ & $-0.016$ \\
\bottomrule
\end{tabular}
\end{table}

The mapping confirms the model's core prediction: \emph{the sign of
the entropy--utility correlation is determined by which uncertainty
source dominates}, not by properties of the entropy signal itself.
Multi-backbone experiments further show that different models shift
the effective balance between Type~I and Type~D uncertainty in the
same environment, making direction unpredictable from environment
structure alone---reinforcing why adaptive direction discovery
(rather than environment-level heuristics) is essential.

\paragraph{Theoretical grounding.}
This direction reversal is an instance of Simpson's
paradox~\citep{simpson1951interpretation, pearl2014understanding}:
aggregating heterogeneous subpopulations with opposing within-group
trends can reverse the aggregate trend. Our two-source decomposition
draws on the epistemic/aleatoric
distinction~\citep{hullermeier2021aleatoric, der2009aleatory},
adapted to the meta-decision setting: Type~I states exhibit
epistemic-like uncertainty (information deficit), while Type~D states
exhibit aleatoric-like diversity (multiple viable paths).
Recent work on LLM calibration~\citep{tao2025revisiting} independently
finds that uncertainty estimation quality varies substantially across
task types---reasoning-oriented tasks yield better-calibrated
estimates than knowledge-seeking ones---providing independent evidence
that uncertainty semantics are task-dependent, not universal.
The key insight is that \emph{the same entropy value has opposite
causal effects on optimizer utility depending on which source
dominates}---a distinction that all prior adaptive compute methods
implicitly collapse.

\begin{proposition}[Necessity of Direction Discovery]
\label{prop:necessity}
Suppose there exist two environments $\mathcal{E}_1, \mathcal{E}_2$
with opposite true directions: $d^*(\mathcal{E}_1) = +1$ and
$d^*(\mathcal{E}_2) = -1$. Then no fixed-direction gate $g_d$ can
simultaneously satisfy
$\mathrm{SR}(g_d, \mathcal{E}_1) \geq \mathrm{SR}(\mathrm{base},
\mathcal{E}_1)$ and
$\mathrm{SR}(g_d, \mathcal{E}_2) \geq \mathrm{SR}(\mathrm{base},
\mathcal{E}_2)$.
That is, direction discovery is a \emph{necessary} condition for
cross-environment non-negative value of computation.
\end{proposition}

\begin{proof}[Proof sketch]
A fixed-direction gate $g_d$ triggers when $d \cdot \sigma(s) > \theta$
for some threshold $\theta$. Consider two cases:

\emph{Case 1:} $d = +1$. The gate triggers on high-signal states
($\sigma(s) > \theta$). In $\mathcal{E}_2$ where $d^* = -1$, high
signal indicates Type~I states where rollouts are harmful
(Eq.~\eqref{eq:direction} with $p_I > p_I^*$). The gate thus
systematically invokes the optimizer at states where
$\mathbb{E}[U] < 0$ and abstains at states where
$\mathbb{E}[U] > 0$, yielding
$\mathrm{SR}(g_{+1}, \mathcal{E}_2) < \mathrm{SR}(\mathrm{base},
\mathcal{E}_2)$.

\emph{Case 2:} $d = -1$. By symmetric argument, the gate fails in
$\mathcal{E}_1$.

Since neither $d = +1$ nor $d = -1$ can satisfy both environments
simultaneously, any gate achieving non-negative VOC in both must
adapt its direction---i.e., direction discovery is necessary.
Empirically, the damage is severe: wrong-direction gating degrades SR
by $-$37.0\,pp on HotpotQA and $-$23.2\,pp on WebShop
(\S\ref{sec:ablation}). Full proof in Appendix~\ref{app:proofs}.
\end{proof}

The model yields three testable predictions:
\begin{enumerate}[nosep,label=\textbf{P\arabic*:}]
\item \textbf{Temporal dynamics.} Within the same environment,
  $\rho$ should \emph{decrease} (become more negative) over the
  episode: early steps mix Type~I and Type~D uncertainty, while
  late steps isolate the residual Type~I component.
\item \textbf{Cross-environment consistency.} Environments with
  similar task structure should exhibit similar $\rho$
  (e.g., FEVER $\approx$ HotpotQA, both search-based QA).
\item \textbf{Signal identity alignment.} In Type~I environments,
  the strongest signal should measure information sufficiency; in
  Type~D environments, decision complexity.
\end{enumerate}
We verify P1--P3 in \S\ref{sec:theory-verification}.


\subsection{Design Implications}
\label{sec:design-implications}

\begin{table}[t]
\caption{Method classification by adaptive compute components. All
prior methods share \{Single signal, Fixed direction, Problem-level\}.
EAAG is the only method with \{Multi-signal, Learned direction,
Step-level, Environment-aware\}.}
\label{tab:classification}
\centering\small
\begin{tabular}{lcccc}
\toprule
Method & Signal & Direction & Gran. & Env. \\
\midrule
CaTS$^\dagger$     & Single & Fixed & Prob. & \xmark \\
SEAG$^\dagger$     & Single & Fixed & Prob. & \xmark \\
CoRefine$^\dagger$ & Single & Fixed & Prob. & \xmark \\
CATTS              & Single & Fixed & Prob. & \xmark \\
AUQ                & Single & Fixed & Prob. & \xmark \\
\midrule
SCG$^\dagger$ (abl.)& Multi & Learned & Step & \cmark \\
\textbf{EAAG}      & \textbf{Multi} & \textbf{Learned}
                    & \textbf{Step} & \cmark \\
\bottomrule
\end{tabular}
\end{table}

Propositions and observations jointly establish that any method
fixing signal, direction, or relying on a single scalar cannot
guarantee non-negative VOC across environments. Direction must be
\emph{discovered}, not assumed. This motivates EAAG (\S\ref{sec:method}).


% ============================================================
% 4. METHOD: EAAG
% ============================================================
\section{Method: Environment-Aware Adaptive Gating}
\label{sec:method}

The analysis in \S\ref{sec:signal-landscape} establishes three
requirements for cross-environment adaptive compute: the gate must
use \emph{multiple signals} (Obs.~2--3), discover the
\emph{direction} of each signal (Prop.~\ref{prop:necessity}), and
adapt to each \emph{environment's} information structure
(Obs.~1, Table~\ref{tab:env-type-mapping}).
We present EAAG (Environment-Aware Adaptive Gating), which satisfies
all three. Unlike prior methods
that assume which signal matters and in which direction, EAAG discovers
both from online interaction data.

\paragraph{Design principles.}
Each component of EAAG follows directly from the analysis in
\S\ref{sec:signal-landscape}:
\begin{enumerate}[nosep,leftmargin=*]
\item \emph{Explore before commit.} Since signal semantics are
  environment-dependent (Obs.~1) and cannot be known a priori, the
  agent must first observe signal--utility relationships in each new
  environment before committing to a gating policy. This motivates
  the exploration phase.
\item \emph{Discover direction, not just threshold.} Fixed-direction
  gates cannot guarantee non-negative value of computation
  (Prop.~\ref{prop:necessity}). The learning phase must recover both
  the \emph{sign} and \emph{identity} of informative signals---a
  requirement naturally satisfied by signed linear weights.
\item \emph{Multi-signal, sparse selection.} Single-signal gates
  carry near-zero information (Obs.~3, AUC${\approx}$0.50). The
  gate must combine multiple signals, but the relevant subset varies
  per environment (Obs.~2). LASSO provides exactly this: automatic
  selection of a sparse, environment-specific feature subset.
\end{enumerate}

\paragraph{Why the method is intentionally simple.}
Our analysis reveals that \emph{the bottleneck for adaptive compute
is not gate complexity but direction discovery}. An MLP gate with
the wrong direction degrades SR by 37.0\,pp
(\S\ref{sec:ablation}); a logistic gate with the right direction
Pareto-dominates all baselines. The complexity budget should go to
\emph{discovery}---identifying which signals matter and in which
direction---not to the gate function itself.

\subsection{Overview}

EAAG operates in three steps (Figure~\ref{fig:method}):
\textbf{Explore} $\to$ \textbf{Reason} $\to$
\textbf{Learn \& Deploy}, with optional online adaptation.

\subsection{Step 1: Explore}

The agent explores a new environment with $N_{\text{explore}} = 50$
episodes using randomized gating ($\varepsilon = 0.5$): at each step,
the optimizer $T$ is invoked with probability 0.5 regardless of state.
This produces a dataset $\mathcal{D} = \{(s_t^{(i)},\; U_t^{(i)})\}$
of (trajectory state, rollout utility) pairs. The randomized design
ensures unbiased estimation of $U$ across the full signal range---a
requirement that targeted exploration (e.g., ``trigger only when
uncertain'') would violate by conditioning on the signal whose
semantics are unknown.

\subsection{Step 2: Reason}

An LLM analyzes $\mathcal{D}$ to identify patterns distinguishing
useful from wasteful rollouts. It outputs:
\begin{itemize}[nosep,leftmargin=*]
\item An \emph{environment profile}: a natural-language description
  of the task structure (e.g., ``early search steps are critical for
  this QA task; evidence availability determines rollout value'').
\item \emph{Feature hypotheses}: task-specific signals to extract
  (e.g., \texttt{price\_mentioned}, \texttt{action\_is\_click} for
  web shopping).
\end{itemize}
These LLM-generated features complement a set of universal features
(step count, token entropy, state length, etc.) available across all
environments.

\subsection{Step 3: Learn \& Deploy}

LASSO selects the most relevant features from the union of universal
and LLM-generated candidates, then trains a logistic regression gate:
\begin{equation}
  g(s) = \mathbf{1}\!\left[\sigma(\mathbf{w}^\top \phi(s) + b)
  > \tau\right]
  \label{eq:gate}
\end{equation}
where $\phi(s)$ is the selected feature vector, $\mathbf{w}$ encodes
both direction (sign) and importance (magnitude) for each signal,
and $\tau = 0.5$ by default. The sign of each $w_i$ directly
implements direction discovery: $w_i > 0$ means ``trigger when signal
$i$ is high'' (Type~D semantics), $w_i < 0$ means ``trigger when
signal $i$ is low'' (Type~I semantics), and $w_i = 0$ means ``this
signal is uninformative in this environment'' (LASSO selection).
Training completes in ${<}1$\,s on CPU.
At deployment, the gate adds zero per-step cost (a single inner
product).

\subsection{Online Adaptation (Optional)}

During deployment, EAAG continues collecting data via
$\varepsilon$-greedy exploration ($\varepsilon: 0.1 \to 0$) and
retrains the gate every 30 episodes, enabling adaptation to
distributional shift.

\paragraph{LLM as environment reasoner: a concrete example.}
The LLM's primary value is \emph{generalizability}, not marginal
accuracy. In most environments, universal features suffice (removing
the LLM changes SR by ${<}$1\,pp). However, in environments with
strong task-specific signals, LLM-generated features are critical.
Consider WebShop: token entropy carries near-zero signal
($\rho{=}{-}0.019$), and no universal feature captures the
environment's decision structure. Given exploration data, the LLM
reasons: ``\emph{The agent struggles most when choosing among
visually similar products. Rollouts help when the agent has found
relevant items but not when it is still searching.}'' From this, it
proposes \texttt{price\_mentioned} (product comparison phase) and
\texttt{action\_is\_click} (commitment vs.\ browsing). LASSO selects
both ($w > 0$), correctly learning that WebShop is Type~D:
high-optionality states benefit from rollouts. Without the LLM, the
gate would rely on step count alone---a weak proxy that misses the
environment's product-comparison structure.


% ============================================================
% 5. EXPERIMENTS
% ============================================================
\section{Experiments}
\label{sec:experiments}

\subsection{Experimental Setup}
\label{sec:setup}

We evaluate EAAG along three axes: (1)~Does environment-aware gating
Pareto-dominate fixed-direction baselines across diverse environments?
(\S\ref{sec:main-results}) (2)~Which components---direction learning,
multi-signal features, LLM reasoning---drive the gains?
(\S\ref{sec:ablation}) (3)~Do the Two-Source Model's predictions
hold empirically? (\S\ref{sec:theory-verification})

\paragraph{Environments.} We evaluate on 8 environments spanning QA,
code generation, web navigation, fact verification, text games, and
manufacturing planning (Table~\ref{tab:env-setup}): HotpotQA,
APPS Intro, WebShop, FEVER, TWExpress, Plancraft, APPS Interview,
and CRUXEval. These environments are selected to cover the full
range of the Two-Source Model---from pure Type~I (FEVER) through
mixed ($\approx p_I^*$, APPS Intro) to Type~D (APPS Interview),
including extreme rollout properties (TWExpress: rollout-safe;
Plancraft: rollout-harmful).

\paragraph{Baselines.} All methods share the same agent and
optimizer $T$; we compare only the gate decision.
(1)~\emph{Bounds}: base\_only, always\_trigger, oracle.
(2)~\emph{Fixed-direction CB}: CaTS, SEAG, CoRefine, CATTS, AUQ,
s1\_budget---all assume high uncertainty $\to$ trigger.
(3)~\emph{Ablation}: SCG (hand-crafted features, requires Phase~1
data), BSW (intentionally reversed direction).
(4)~\emph{Ours}: EAAG.

\paragraph{Metrics.} SR (success rate, $\uparrow$) and Cost
(rollouts per episode during deployment, $\downarrow$).
A method Pareto-dominates another if $\text{SR} \geq$ and
$\text{Cost} \leq$ with at least one strict inequality.
Cost measures only the \emph{deployment} rollout budget---calibration
or exploration overhead is excluded for all methods, ensuring a fair
comparison of runtime efficiency.

\begin{table}[t]
\caption{Environment setup. Base/Always: SR without/with optimizer.
$\Delta$: rollout headroom.}
\label{tab:env-setup}
\centering\small
\begin{tabular}{lcccl}
\toprule
Env. & Base & Always & $\Delta$ & Optimizer $T$ \\
\midrule
HotpotQA      & 49.0 & 97.0 & +48.0 & Per-action eval \\
APPS Intro    & 58.5 & 64.5 &  +6.0 & $K$-variant \\
WebShop       &  7.2 & 43.0 & +35.8 & LLM-Propose-$K$ \\
FEVER         & 37.0 & 99.8 & +62.8 & Per-action eval \\
TWExpress     & 67.5 & 99.3 & +31.8 & Per-action eval \\
Plancraft     & 29.8 & 22.8 &  $-$7.0 & Per-action eval \\
APPS Intv.    & 60.5 & 79.5 & +19.0 & $K$-variant \\
CRUXEval      & 85.0 & 99.5 & +14.5 & $K$-variant \\
\bottomrule
\end{tabular}
\end{table}


\subsection{Main Results}
\label{sec:main-results}

Table~\ref{tab:main} reports SR and deployment Cost across 8
environments. EAAG achieves 34 wins (18 statistically significant
at 95\% bootstrap CI) against 2 losses in head-to-head
SR comparisons with 6 baselines across 8 environments
(Table~\ref{tab:winloss}). Both losses are on HotpotQA against
AUQ ($-$1.8\,pp, CI $[-4.0, +0.3]$) and s1\_budget ($-$1.8\,pp,
CI $[-4.2, +0.3]$)---neither is statistically significant (CIs
include zero), and both occur in a ceiling-effect environment
where EAAG (95.2\%) already approaches oracle (97.0\%).
EAAG Pareto-dominates CaTS in 5/6
environments (the exception is FEVER, where EAAG's 49.8\% trails
CaTS's 50.2\% by 0.4\,pp---within noise---at less than half the
cost) and SCG in 4/7 environments (the latter requiring Phase~1
data that EAAG eliminates).

\paragraph{The story of the results.}
Three patterns stand out. \emph{First}, the gap between EAAG and
fixed-direction baselines is \emph{largest} precisely where the
Two-Source Model predicts direction mismatch. In FEVER (Type~I
dominated, $\rho_{\mathrm{entropy}}{=}{-}0.119$), all ``high
uncertainty $\to$ trigger'' baselines cluster around 50\% while the
oracle reaches 99.8\%; the assumed direction is wrong, and no amount
of calibration rescues them. \emph{Second}, EAAG's advantage is
\emph{smallest} where direction matters least: HotpotQA approaches a
ceiling (95.2\% vs.\ oracle 97.0\%) because the massive rollout
headroom (+48\,pp) means even imprecise gating works.
\emph{Third}, EAAG's behavior is \emph{qualitatively different} from
baselines: it does not simply improve numbers but exhibits
environment-adaptive strategies---aggressive triggering in
high-headroom environments, conservative gating in low-headroom
ones, and near-zero triggering when rollouts are harmful.

% TODO: Insert main results table (tab:main) with full method x environment SR/Cost grid
% TODO: Insert win/loss summary table (tab:winloss)

\paragraph{Per-environment analysis.}
\begin{itemize}[nosep,leftmargin=*]
\item \textbf{WebShop} (strongest differentiation):
  EAAG achieves 43.8\% SR ($\approx$ oracle
  43.0\%) at 2.29 ro/ep---vs.\ CaTS 30.5\% at 3.05 ro/ep. The gate
  triggers selectively (16.9\% of steps) with 75.1\% precision.
  This is where LLM-generated features
  (\texttt{price\_mentioned}, \texttt{action\_is\_click}) contribute
  most, compensating for entropy's near-zero signal
  ($\rho{=}{-}0.019$).
\item \textbf{FEVER} (direction mismatch exposed):
  All fixed-direction CBs cluster around 50\%
  (CaTS 50.2\%, SEAG 49.3\%, CoRefine 49.8\%), far below
  always\_trigger (99.8\%). CATTS fails at 34.2\% $<$ base 37.0\%.
  EAAG achieves 49.8\%, limited by online exploration bias
  (\S\ref{sec:discussion})---an honest limitation that itself
  validates the theory (critical rollouts concentrate at step~0,
  where random exploration misses 50\% of the time).
\item \textbf{HotpotQA} (ceiling effect):
  Most methods approach oracle (97.0\%). EAAG achieves 95.2\% at
  lowest cost (1.34 ro/ep). The advantage here is efficiency, not
  accuracy: EAAG uses 38\% fewer rollouts than the next-best method.
\item \textbf{APPS Intro} (narrow headroom):
  Only +6\,pp headroom. EAAG learns moderate gating
  (RR=35\%), achieving 66.0\% at 1.20 ro/ep---the lowest cost
  among methods that improve over base.
\end{itemize}

\subsection{Environments with Extreme Rollout Properties}
\label{sec:extreme}

TWExpress and Plancraft represent extreme ends of the rollout value
spectrum, testing whether EAAG adapts beyond ``normal'' operating
ranges.

TWExpress (rollout-safe, $\Delta{=}+$31.8\,pp): EAAG achieves
99.0\% (vs.\ always 99.3\%), learning to trigger aggressively
(RR=73\%)---correctly inferring that rollouts are almost
always beneficial.
Plancraft (rollout-harmful, $\Delta{=}-$7.0\,pp): EAAG achieves
23.3\% (vs.\ always 22.8\%), with a trigger rate that decays over
steps (49\% at step~0 $\to$ $<$20\% at later steps), learning to
become increasingly conservative as evidence of rollout harm
accumulates.


\subsection{Ablation Studies}
\label{sec:ablation}

The main results show \emph{that} EAAG works; we now ask
\emph{why}. We ablate three components in decreasing order of
expected importance: direction learning, multi-signal features, and
LLM reasoning.

\paragraph{Direction is the primary determinant (BSW ablation).}
Intentionally reversing the learned direction causes catastrophic
failure across environments: BSW achieves 58.2\% vs.\ EAAG's 95.2\%
on HotpotQA ($-$37.0\,pp) and 20.6\% vs.\ 43.8\% on WebShop
($-$23.2\,pp).
On HotpotQA, the MLP gate with wrong direction achieves 45.3\%---
\emph{below} the 49.0\% no-trigger baseline. This confirms the
central prediction of Proposition~\ref{prop:necessity}: direction is
not a minor calibration detail but the primary determinant of gate
quality. The magnitude of degradation correlates with the
environment's signal strength ($|\rho|$ vs.\ $|\Delta\mathrm{SR}|$,
$R^2 > 0.5$; Appendix~\ref{app:proofs}), as the Two-Source Model
predicts.

\paragraph{LLM reasoning provides robustness, not accuracy.}
Removing the LLM reasoning step (using only universal features)
changes SR by ${<}$1\,pp in 4 of 6 environments: 95.2\%$\to$95.8\%
(HotpotQA), 66.0\%$\to$65.8\% (APPS Intro), 43.8\%$\to$43.7\% (WebShop),
99.0\%$\to$99.2\% (TWExpress).
Two environments show larger effects: FEVER (+9.1\,pp with LLM,
where LLM-generated features capture step-0 criticality) and
Plancraft ($-$3.9\,pp, where LLM features introduce noise in a
rollout-harmful environment). The LLM's value is not
marginal SR improvement but \emph{generalizability}: it enables
zero-shot feature generation for unseen environments where
task-specific signals (e.g., WebShop's \texttt{price\_mentioned})
cannot be anticipated by universal features alone.

\paragraph{Gating magnitude emerges from direction learning.}
The learned gate adapts trigger rate across environments:
aggressive triggering in
rollout-safe environments (TWExpress: RR=73\%, HotpotQA: RR=66\%),
moderate in mixed environments (APPS Intro: RR=35\%, WebShop: RR=28\%),
and step-dependent in rollout-harmful ones (Plancraft: starts at
49\% but decays to $<$20\% at later steps as the gate learns rollouts
are counterproductive). These environment-adaptive gating patterns
arise from simple logistic regression without explicit headroom
estimation---the
LASSO coefficients implicitly encode environment-specific triggering
patterns through the learned signal--utility relationship.


\subsection{Two-Source Model Verification}
\label{sec:theory-verification}

The ablations above confirm EAAG's practical advantage; we now test
whether the \emph{theoretical mechanism}---the Two-Source Model
(\S\ref{sec:toy-model})---correctly predicts the observed patterns.
Following \citet{schaeffer2023emergent}, we derive testable
predictions from the model and verify each against held-out data.

\paragraph{P1: Temporal dynamics (confirmed).}
\emph{Prediction}: $\rho$ should
\emph{decrease} (become more negative) over the episode, because
early steps mix Type~I and Type~D uncertainty while late
steps isolate the residual Type~I component.
\emph{Test}: We split trajectory data into early
(step $\leq$ median) and late (step $>$ median) subsets and compute
$\rho(\text{entropy}, U)$ separately.
\emph{Result}: Confirmed across all environments with sufficient
signal. In HotpotQA, $\rho$ shifts from $-0.176$ (early,
$p{<}10^{-4}$) to $-0.437$ (late, $p{<}10^{-5}$): late-step
entropy reflects \emph{deeper} information poverty after failed
evidence retrieval. In code generation (APPS Intro), $\rho$ shifts from
$+0.102$ (early, non-significant) to $-0.144$ (late, $p{=}0.04$):
early exploration diversity gives way to debugging frustration.
In WebShop, $\rho$ shifts from $+0.285$ (early, $p{<}10^{-5}$)
to $-0.006$ (late, non-significant): the strong Type~D signal
during product browsing vanishes once the agent commits.

\paragraph{P2: Cross-environment consistency (confirmed).}
\emph{Prediction}: Environments with similar information structure
should exhibit similar $\rho$ signs and magnitudes.
\emph{Test}: We compare same-family pairs.
\emph{Result}: FEVER ($\rho{=}{-}0.119$) and HotpotQA
($\rho{=}{-}0.041$), both search-based QA tasks requiring evidence
retrieval, share negative entropy--utility direction.
APPS Intro ($\rho{=}{+}0.012$) and APPS Interview
($\rho{=}{+}0.317$), both code generation tasks, share non-negative
direction. The magnitude difference within each pair reflects the
degree of Type~I contamination.

\paragraph{P3: Signal identity alignment (confirmed).}
\emph{Prediction}: In Type~I environments, the strongest signal
should measure information sufficiency; in Type~D/mixed environments,
decision complexity.
\emph{Test}: We examine the LASSO-selected features per environment.
\emph{Result}: In Type~I environments (HotpotQA, FEVER), the
strongest signal (\texttt{step\_count}, $|\rho|{=}0.494$ and $0.619$)
measures trajectory progress---a proxy for information accumulation.
In WebShop (mixed, choice-dominated),
\texttt{num\_available\_actions} ($\rho{=}{+}0.444$) measures the
decision space size---a proxy for decision complexity. All three
predictions are confirmed, providing converging evidence for the
Two-Source Model.

\paragraph{Controlled information manipulation (supplementary).}
To provide stronger evidence that information structure drives signal
semantics, we manipulate information availability within HotpotQA
itself. \emph{InfoPoor} (search returns only the first sentence)
reduces base SR to 44.0\% while \emph{InfoRich} (gold evidence
injected at start) raises it to 82.0\%. The signal hierarchy shifts:
in InfoPoor, \texttt{step\_count} dominates ($\rho{=}{-}0.608$) while
entropy carries a weak positive signal ($\rho{=}{+}0.119$);
in InfoRich, entropy becomes the dominant positive
signal ($\rho{=}{+}0.311$) while \texttt{step\_count} weakens
($\rho{=}{-}0.147$). This demonstrates that the \emph{identity} of
the most informative signal---not just its direction---is a function
of the environment's information structure.


\subsection{Robustness of Direction Reversal}
\label{sec:robustness}

A natural concern is whether the observed direction reversal is a
genuine phenomenon or an artifact of confounders. We address this
through three complementary analyses.

\paragraph{Stratified analysis: controlling for trajectory length.}
Trajectory length varies across environments and correlates with both
entropy and utility. To rule out length as a confounder, we stratify
by step count and recompute $\rho(\text{entropy}, U)$ within each
stratum.
\emph{Result}: Direction reversal persists within trajectory-length
strata where utility variance is non-zero. In HotpotQA, $\rho$ remains
negative across all three strata ($-0.18$/$-0.46$/$-0.42$).
The reversal is not an artifact of trajectory length.

\paragraph{Interventional evidence: wrong-direction ablation.}
The BSW ablation (\S\ref{sec:ablation}) provides \emph{interventional}
evidence: we \emph{manipulate} the gate's direction and measure the
causal effect on SR. If direction reversal were merely correlational,
reversing direction should have no systematic effect. Instead, we
observe catastrophic degradation:
$-$37.0\,pp on HotpotQA ($d^*{=}{-}1 \to d{=}{+}1$) and
$-$23.2\,pp on WebShop. The MLP gate with wrong direction falls
\emph{below} the no-trigger baseline (45.3\% $<$ 49.0\%), confirming
that the direction encodes a real causal relationship, not a
statistical artifact.

\paragraph{Gate complexity ablation: direction $\gg$ capacity.}
To verify that the bottleneck is direction rather than gate
complexity, we compare gate architectures with correct vs.\ wrong
direction (Table~\ref{tab:capacity}).

\begin{table}[t]
\caption{Gate complexity ablation on HotpotQA. Direction matters
more than model capacity: a logistic gate with correct direction
outperforms an MLP with wrong direction by 49.9\,pp.}
\label{tab:capacity}
\centering\small
\begin{tabular}{lccc}
\toprule
Gate & Direction & SR (\%) & $\Delta$ \\
\midrule
Logistic (5 feat) & Correct & 95.2 & +46.2 \\
MLP (5 feat)      & Correct & $\sim$95 & $\sim$+46 \\
Hidden state (2560-d) & Correct & $\sim$95 & $\sim$+46 \\
\midrule
Logistic (5 feat) & Wrong & 62.0 & +13.0 \\
MLP (5 feat)      & Wrong & 45.3 & $-$3.7 \\
\midrule
No gate (base)    & --- & 49.0 & 0 \\
\bottomrule
\end{tabular}
\end{table}

With correct direction, all architectures achieve $\sim$95\% SR;
the jump from logistic to MLP to hidden-state probe adds $<$1\,pp.
With wrong direction, higher capacity makes performance \emph{worse}
(MLP 45.3\% $<$ logistic 62.0\%), because the more powerful gate more
precisely targets harmful states. The information hierarchy is:
\textbf{signal direction $\gg$ signal count $\gg$ gate complexity}.


% ============================================================
% 6. DISCUSSION
% ============================================================
\section{Discussion}
\label{sec:discussion}

\paragraph{Insight for the community.}
Our central message is that \emph{the bottleneck for adaptive compute
is the understanding of uncertainty semantics}. The same entropy value
means ``the agent lacks information'' in one environment and ``the
agent faces multiple viable paths'' in another. Any method that
collapses this distinction will systematically fail in at least one
environment type (Proposition~\ref{prop:necessity}). We offer a
prescriptive framework: before designing a gating mechanism for a
new environment, first characterize its information structure
(information sufficiency, decision reversibility, feedback delay) to
predict signal semantics and direction.

\paragraph{FEVER and exploration bias.}
EAAG achieves 49.8\% on FEVER---far below SCG's 98.0\% (which uses
Phase~1 always-trigger data) and always-trigger's 99.8\%. Analysis
reveals that FEVER's rollout value concentrates in step~0--1
($\rho{=}{-}0.619$). Random exploration misses step~0 with 50\%
probability, producing late-step negative-utility data that biases
LASSO toward never triggering. This is an inherent limitation of
online exploration in environments where critical decisions occur
in the first step. Notably, this failure mode is itself predicted by
the Two-Source Model: FEVER is the most extreme Type~I environment
($p_I$ near 1.0), where nearly all uncertainty is informational---
precisely the regime where undirected exploration is least efficient.

\subsection{Future Directions}

\paragraph{Beyond two sources: a taxonomy of uncertainty types.}
Our Two-Source Model (Type~I / Type~D) is a first-order
approximation. Real environments likely contain finer-grained
uncertainty types---e.g., \emph{execution uncertainty} (the agent
knows what to do but may fail mechanically), \emph{compositional
uncertainty} (sub-goals are clear but their interaction is not), and
\emph{temporal uncertainty} (information will arrive but has not yet).
Each type implies a different relationship between uncertainty signals
and optimizer utility. A richer taxonomy, grounded in the structure of
environment MDPs, would enable more precise signal--utility modeling.

\paragraph{Adaptive exploration for step-0 critical environments.}
EAAG's FEVER limitation reveals a concrete open problem:
\emph{how should an agent explore when the critical decision window
is narrow?} Uniform random exploration wastes half its budget on
non-critical steps. Curriculum-based exploration---starting with
always-trigger and gradually introducing gating---could address this.
More ambitiously, a meta-learner could transfer structural knowledge
across environments to warm-start exploration.

\subsection{Limitations}

\begin{enumerate}[nosep,leftmargin=*]
\item \textbf{Backbone diversity.} We verify our core findings on
  three backbones (Qwen3-4B, Phi-3.5-mini-instruct, Llama-3.1-8B-Instruct)
  across all 8 environments (Appendix~\ref{app:multi-backbone}).
  The multi-backbone results reveal that signal-utility direction depends on
  \textbf{both} the environment \emph{and} the model backbone.
  Among 5 environments with significant entropy signals on
  both new models, 4/5 show cross-model direction variation.
  This reinforces the need for adaptive gating.
  Testing on larger models (70B+) remains future work.
\item \textbf{Linear Two-Source Model.} The model assumes linear
  conditional relationships ($U_I \sim -\alpha H$,
  $U_D \sim +\beta H$). Real environments may exhibit nonlinear
  mixtures or more than two uncertainty sources. The model's value
  is explanatory and prescriptive, not predictive of exact $\rho$ values.
\item \textbf{$p_I$ estimation.} The Two-Source Model uses $p_I$
  (fraction of Type~I states) as a theoretical construct. Estimating
  $p_I$ from data without ground-truth state labels remains open.
  EAAG sidesteps this by learning direction \emph{end-to-end}
  without explicit $p_I$ estimation.
\item \textbf{Exploration cost.} EAAG requires 50 exploration episodes
  per new environment. While modest relative to thousands of
  deployment episodes, this is non-zero. The cost is amortized:
  once direction is learned, no further exploration is needed unless
  the environment's information structure shifts.
\end{enumerate}

\subsection{Broader Impact}

This work addresses computational efficiency in LLM agent deployment.
By reducing unnecessary optimizer invocations, EAAG lowers the energy
and cost footprint of test-time compute without sacrificing task
performance. The prescriptive framework (characterize information
structure before designing gates) may help practitioners avoid the
pitfall of deploying fixed-direction methods that actively harm
performance in mismatched environments. We do not foresee direct
negative societal impacts beyond those inherent to the LLM agents
themselves.


% ============================================================
% 7. CONCLUSION
% ============================================================
\section{Conclusion}

We identify a hidden assumption shared by all existing adaptive
test-time compute methods: that the mapping from uncertainty signals
to compute need has a fixed direction. Through systematic measurement
across 8 diverse agent environments and 3 model backbones, we show
that this assumption is wrong---signal--utility direction reverses
across environments \emph{and} across model backbones, making it
fundamentally unpredictable without adaptive learning. We
formalize this via a Two-Source uncertainty model that explains when
and why direction reverses, and prove that direction discovery is a
necessary condition for cross-environment non-negative value of
computation. EAAG, which lets the LLM agent autonomously discover
environment-specific gating patterns through exploration, reasoning,
and sparse direction learning, achieves 34 wins (18 statistically
significant) vs.\ 2 losses against 6 competing methods across 8
environments and exhibits environment-adaptive gating patterns
without explicit headroom estimation.
The bottleneck was never the method's complexity---it was the
assumption.


% ============================================================
% REFERENCES
% ============================================================
\bibliography{references}
\bibliographystyle{acl_natbib}


% ============================================================
% APPENDIX
% ============================================================
\appendix

\section{Extended Related Work}
\label{app:related}

% Detailed method-by-method comparison table
% Timeline of concurrent submissions
% Full landscape of adaptive compute methods

\paragraph{Detailed method comparison.}
\begin{itemize}[nosep,leftmargin=*]
\item \textbf{CATTS}~\citep{catts2026}: Vote entropy + margin; $K{=}5$
  forward passes per step; direction = ``high disagreement $\to$
  trigger''; overhead = $5\times$ per step.
\item \textbf{SEAG}~\citep{seag2025}: Mean token confidence; Platt
  scaling; problem-level; direction = ``low confidence $\to$ search
  deeper.''
\item \textbf{CaTS}~\citep{cats2025}: Calibrated confidence +
  Best-of-N early stopping; direction = ``low confidence $\to$
  generate more candidates.''
\item \textbf{CoRefine}~\citep{corefine2026}: Conv1D confidence
  controller; halt/re-examine/redirect; direction = ``low confidence
  $\to$ refine.''
\item \textbf{AdaptThink}~\citep{adaptthink2025}: RL-learned
  think/no-think; implicitly learns direction but per-environment
  training, no interpretable signal.
\item \textbf{DiffAdapt}~\citep{diffadapt2025}: U-shaped entropy
  pattern + hidden-state probe; assumes universal entropy shape
  across tasks.
\end{itemize}

All share single-signal or fixed-direction assumptions.


\section{Experimental Details}
\label{app:experiments}

\subsection{Environment Specifications}

For each of the 8 environments, we provide:
\begin{itemize}[nosep,leftmargin=*]
\item Task description and action space
\item State representation and reward function
\item Episode length distribution
\item Optimizer $T$ implementation details:
  \begin{itemize}[nosep]
  \item HotpotQA, FEVER, TWExpress, Plancraft: per-action evaluation
    (generate $K{=}5$ candidates, evaluate each, select best)
  \item APPS Intro, APPS Interview, CRUXEval: $K$-variant sampling
    (generate $K{=}3$ code solutions, test each, select passing one)
  \item WebShop: LLM-Propose-$K$ (LLM proposes $K{=}3$ action
    sequences, simulate each, select highest-reward)
  \end{itemize}
\end{itemize}

\subsection{Hyperparameter Configuration}

\paragraph{EAAG hyperparameters.}
$N_{\text{explore}} = 50$ episodes,
$\varepsilon_{\text{explore}} = 0.5$,
$\varepsilon_{\text{deploy}} = 0.1 \to 0$ over 100 episodes,
LASSO $\alpha$ selected via 5-fold CV,
gate threshold $\tau = 0.5$,
retrain interval = 30 episodes.

\paragraph{Baseline hyperparameters.}
CaTS: Platt scaling on Phase~1 data (200 episodes), threshold = 0.5.
SEAG: mean confidence threshold calibrated on Phase~1.
CoRefine: entropy threshold calibrated on Phase~1.
CATTS: $K{=}5$ votes, entropy threshold = 0.8, margin threshold = 0.2.
SCG: hand-crafted features, direction from Phase~1 data.
BSW: EAAG with intentionally reversed direction (sign flip).

\subsection{Feature Selection Details}

% Table: Environment x Feature matrix showing LASSO coefficients
% LLM-generated features per environment with prompts and outputs

\subsection{Computational Cost Breakdown}

% Table: Method x Cost component
% EAAG: 50 exploration, ~2 GPU-hrs, 0ms gate, <1s train
% CaTS: 200 always-trigger, ~8 GPU-hrs, 0ms gate
% CATTS: 0, 0, K x forward_pass ms per step


\section{Proofs}
\label{app:proofs}

\subsection{Proof of Proposition~\ref{prop:necessity}
  (Necessity of Direction Discovery)}

\begin{proposition*}[Restated]
Suppose there exist two environments $\mathcal{E}_1, \mathcal{E}_2$
with opposite true directions: $d^*(\mathcal{E}_1) = +1$ and
$d^*(\mathcal{E}_2) = -1$. Then no fixed-direction gate $g_d$ can
simultaneously satisfy
$\mathrm{SR}(g_d, \mathcal{E}_1) \geq \mathrm{SR}(\mathrm{base},
\mathcal{E}_1)$ and
$\mathrm{SR}(g_d, \mathcal{E}_2) \geq \mathrm{SR}(\mathrm{base},
\mathcal{E}_2)$.
\end{proposition*}

\begin{proof}
Define a fixed-direction gate as $g_d(s) = \mathbf{1}[d \cdot \sigma(s) > \theta]$
for $d \in \{+1, -1\}$ and threshold $\theta$.

Define $d^*(\mathcal{E}) = +1$ if $\mathrm{Corr}(\sigma(s), U(T,s)) > 0$
in $\mathcal{E}$ (Type~D), and $d^*(\mathcal{E}) = -1$ if
$\mathrm{Corr}(\sigma(s), U(T,s)) < 0$ (Type~I).

\textbf{Lemma (Wrong-direction harm).} If $d \neq d^*(\mathcal{E})$, then
$\mathrm{SR}(g_d, \mathcal{E}) \leq \mathrm{SR}(\mathrm{base}, \mathcal{E}) - \delta(\mathcal{E})$
where $\delta(\mathcal{E}) > 0$ depends on signal strength $|\rho(\mathcal{E})|$.

\emph{Proof of Lemma}: By definition of $d^*$, states with high
$d \cdot \sigma(s)$ are states where $U < 0$. The gate triggers $T$
at these states, causing the agent to adopt worse actions. The base
policy, which never triggers $T$, avoids this systematic harm.

\textbf{Main argument:}

\emph{Case 1}: $d = +1$. In $\mathcal{E}_2$ where $d^* = -1$,
$d \neq d^*(\mathcal{E}_2)$. By Lemma:
$\mathrm{SR}(g_{+1}, \mathcal{E}_2) \leq \mathrm{SR}(\mathrm{base}, \mathcal{E}_2) - \delta(\mathcal{E}_2)
< \mathrm{SR}(\mathrm{base}, \mathcal{E}_2)$.
Constraint violated.

\emph{Case 2}: $d = -1$. In $\mathcal{E}_1$ where $d^* = +1$,
by symmetric argument, constraint violated.

Since $d \in \{+1, -1\}$ and both cases lead to violation, no fixed
$d$ satisfies both constraints simultaneously.

\textbf{Empirical grounding:}
$\delta(\text{HotpotQA}) \approx 37.0$\,pp (BSW ablation),
$\delta(\text{WebShop}) \approx 23.2$\,pp,
$\delta(\text{FEVER}) \approx 36.8$\,pp.
MLP with wrong direction: 45.3\% $<$ base 49.0\% on HotpotQA.
\end{proof}

\subsection{VOC Non-Negativity Scope}

In the classical VOC framework~\citep{russell1991right},
$\mathrm{VOC}(T, s) = \mathbb{E}[\max(V(a_T), V(a_{\mathrm{base}}))] - V(a_{\mathrm{base}}) \geq 0$
because the agent can \emph{disregard} the computation result if
unfavorable. Our setting violates this: when the agent uses an
evaluator/optimizer $T$, the evaluator's assessment \emph{is} the
agent's decision. The agent cannot separately evaluate the
evaluator's output because the evaluator \emph{is} the evaluation
mechanism.

Under evaluator-executor identity, VOC can be negative:
$\mathrm{VOC}(T, s) = \mathbb{E}[V(a_T(s))] - V(a_{\mathrm{base}}(s))$.
If $T$ systematically selects worse actions in Type~I states,
$\mathrm{VOC} < 0$ for those states. This explains why wrong-direction
gating is not merely wasteful ($\mathrm{VOC} = 0$) but actively
harmful ($\mathrm{VOC} < 0$).


\section{Two-Source Model: Full Derivation}
\label{app:model}

\subsection{Full Derivation}

At each step $t$, state $s_t$ belongs to one of two latent types:
$s_t \in \text{Type~I}$ with probability $p_I(\mathcal{E})$,
$s_t \in \text{Type~D}$ with probability $1 - p_I(\mathcal{E})$.

Conditional utility models:
\begin{align}
  U(T, s) \mid s \in \text{Type~I} &= -\alpha \cdot H(s) + \varepsilon_I, \quad \alpha > 0 \\
  U(T, s) \mid s \in \text{Type~D} &= +\beta \cdot H(s) + \varepsilon_D, \quad \beta > 0
\end{align}

The marginal correlation:
\begin{align}
  \mathrm{sign}(\rho(\mathcal{E})) &= \mathrm{sign}(\beta \cdot (1 - p_I) - \alpha \cdot p_I) \\
  &= \mathrm{sign}(\beta - (\alpha + \beta) \cdot p_I)
\end{align}

Reversal threshold: $p_I^* = \beta / (\alpha + \beta)$.

\subsection{Environment Mapping on the $p_I$ Axis}

Estimated positions on the $p_I$ axis (0 = pure Type~D, 1 = pure Type~I):
\begin{itemize}[nosep,leftmargin=*]
\item $p_I \approx 0.2$: APPS Interview (strong positive $\rho = +0.317$)
\item $p_I \approx 0.45$: APPS Intro (near-zero $\rho = +0.012$)
\item $p_I \approx 0.5$: WebShop (entropy $\rho \approx 0$, but non-entropy signals work)
\item $p_I \approx 0.55$: CRUXEval (weak negative $\rho = -0.048$)
\item $p_I \approx 0.6$: Plancraft (weak negative, rollouts harmful)
\item $p_I \approx 0.7$: HotpotQA (negative $\rho = -0.041$)
\item $p_I \approx 0.8$: TWExpress (negative $\rho = -0.290$)
\item $p_I \approx 0.9$: FEVER (strong negative $\rho = -0.119$)
\end{itemize}


\section{Multi-Backbone Verification}
\label{app:multi-backbone}

We validate our core findings on three backbones from three vendors:
Qwen3-4B-Instruct (Alibaba, 4B), Phi-3.5-mini-instruct (Microsoft, 3.8B),
and Llama-3.1-8B-Instruct (Meta, 8B) across all 8 environments.

\subsection{Signal Direction: Environment $\times$ Model Interaction}

The multi-backbone results (Table~\ref{tab:signal-discovery} in main text)
reveal that 4/5 environments with $\geq$2 significant signals show
cross-model direction variation. Structural signals
(\texttt{step\_count}) are model-invariant while entropy signals are
model-dependent.

\subsection{Fixed-Direction Baselines Fail Across Backbones}

On Phi-3.5 APPS: all 6 CB methods score \emph{below} base\_only
(59\% $\to$ 27--36\%), demonstrating that fixed-direction assumption
is catastrophically wrong when the model changes.


\section{Additional Analyses}
\label{app:additional}

\subsection{Hyperparameter Sensitivity}

EAAG is robust to $N_{\text{explore}} \in [30, 100]$ (stable SR for
$N_{\text{explore}} \geq 30$; marginal gains above 50) and gate
threshold $\tau \in [0.3, 0.7]$.

\subsection{Statistical Significance}

All head-to-head comparisons use 3-seed evaluation (200 episodes/seed).
Bootstrap confidence intervals (5000 resamples) are reported for all
comparisons. 18 of 34 wins are statistically significant (95\% CI
excludes zero). Pareto-dominance claims use a conservative criterion
that does not rely on p-values.


\end{document}
